\section{Arithmetic estimates for the exact reconstruction}
\label{p2:technical-transfer}

We use the notation of Section~\ref{p13:transfer:section} of the article.
All original windows with $p_x>0$ remain present. In this section,
$\rho$ means the effective relative nullity $\rho_{\mathrm{rel}}$.
The estimates culminate in separate bounds for each fixed connected
order and for the entire remaining range of cardinalities. The article
then assembles them using the positive stopped expansion.

\subsection{Two dense targets above a prescribed space}

The independent counts do not control a large rank gain: a family of
nullity $d$ carries an additional factor $2^d$. We now isolate the
conditional geometric estimate that pays for this factor. Only two
centers remain free. Every prime circuit meeting at least three
physical blocks must then meet one of the fixed sources, which makes
its coefficient enumerable by a divisor bound.

\begin{lemma}[Conditional dense pair estimate]
\label{p13:transfer:dense-targets}
Fix a number $s$ of source windows and an effective subspace $W$ of
their value space. Suppose $S_y\cap W=S_z\cap W=0$, and write
\[
 d=\dim W+\dim S_y+\dim S_z-\dim(W+S_y+S_z).
\]
For every fixed $\eta,\varepsilon>0$ there exists a fixed $\theta>0$
such that, uniformly in the sources and $W$,
\begin{equation}
 \sum_{\substack{y\ne z,\ y,z\text{ outside the sources}\\
             d\geq\eta L\\
             (y,z)\text{ nonaligned or }\dim(S_y\cap S_z)\leq\theta L}}
         p_yp_z2^d\ll_{s,C,\eta,\varepsilon}N^{-1+\varepsilon}.
 \label{p13:transfer:dense-target-bound}
\end{equation}
In particular the hypothesis on the pair is satisfied when $S_y,S_z$
are direct. For a primitive family of fixed size $r\geq3$, with a
prescribed collection of $r-2$ sources having at most $N^{a+\varepsilon}$
possibilities, its dense homogeneous mass is at most
\begin{equation}
 \sum_{\rho\geq\eta L}\Bigl(\prod_{i\in I}p_i\Bigr)2^{\rho(I)}
          \ll N^{-(r-1)+a+\varepsilon}.
 \label{p13:transfer:restricted-dense}
\end{equation}
\end{lemma}

The kernel defining $d$ projects injectively onto each target, so its
coordinate labels force many binary components in each target separately.
After regularization and removal of the weak aligned components, let
$m_t$ count the remaining components meeting target $t$, and let $n_{1,t}$
count those using exactly one of its $B=L+1$ coordinates. Then
$n_{1,t}\geq2m_t-B$. Above $d=(1/2+\delta)L$, this forces a positive
proportion of singleton occurrences in both targets, allowing the stronger
Pell or degree-one curve count. The resulting bounds are
$N^{1/2+\varepsilon'}$ target pairs in the lower stratum and
$N^{\varepsilon'}$ in the upper stratum; multiplying by $p_yp_z2^d$
gives $N^{-1+\varepsilon}$ after choosing the internal margins.
Section~\ref{p14:comp:transfer-dense-targets} supplies the separate
component counts, the coefficient enumeration from fixed sources, and
the removal of small aligned curves, including effective local sections.


\begin{lemma}[A dense compressed target]
\label{p13:transfer:compressed-target}
For fixed sources and a fixed channel of good pairs, the number of
compressed centers $X$ for which
$\dim(S(U_X)\cap W)\geq\eta L$ is $O_{s,C,\eta,\varepsilon}(N^\varepsilon)$.
\end{lemma}

\begin{proof}
The compression has fixed width $O(B)$ and an injective even-selection
map. In Proposition~\ref{p1:fixed:targets}, take $H$ larger than this
width. Its regularization loses a fixed small fraction of $\eta L$.
A prime above $H$ has at most one target occurrence, so all other
occurrences in its circuit are source occurrences. The circuit products
are divisors of fixed source values, and the fixed-source CRT and
split-product steps in that proposition are unchanged. The proof uses
the width bound and the target projection rank, not a target space of
dimension exactly $L$. The finitely many channel dilations and
$B^{O(1)}$ offset patterns are absorbed in the margins.
\end{proof}

\subsection{The exact weights of primitive families}

The primitive estimate below combines the dense-host construction of
Section~\ref{p14:comp:confinement}, the homogeneous envelope of
Lemma~\ref{p2:homogeneous:envelope}, and the conditional half-count
through Corollary~\ref{p13:half:connected}.
We first record a separate small-kernel estimate for related counting
arguments; it is not used in the proof of
Proposition~\ref{p13:transfer:primitive}.
For $Y\geq H$, put
\[
 K_Y(n)=\prod_{\substack{\ell>Y\\v_\ell(n)\ \mathrm{odd}}}\ell,\qquad
 a_Y(n)=\prod_{\ell>Y}\ell^{\lfloor v_\ell(n)/2\rfloor}.
\]
Then $n=s_Y(n)a_Y(n)^2K_Y(n)$, where $s_Y(n)$ is $Y$-smooth.

\begin{lemma}[Several small kernels in one window]
\label{p13:transfer:small-kernels}
Fix $A>2$, $0<b,u<1$, and an integer $m$ with $m(1-b)>1$.
For $Y=(\log N)^A$, the number of windows containing $m$ distinct
values with $K_Y(n)\leq N^b$ and $a_Y(n)^2\leq N^u$ is
\begin{equation}
 \ll N^{\max(1-1/A,(1+u)/2)+\varepsilon}.
 \label{p13:transfer:small-kernels-bound}
\end{equation}
Also, with the previously fixed cutoff $H=C_DB$, the set
$\mathcal X_q$ of windows containing distinct $n_1,n_2$
with $K_H(n_1)K_H(n_2)\leq N^q$ has cardinality
$O(N^{q+\varepsilon})$, for fixed $q\geq0$.
\end{lemma}

\noindent The complete proof is given in Section~\ref{p14:comp:transfer-small-kernels}.


\begin{proposition}[Homogeneous mass at fixed primitive arity]
\label{p13:transfer:primitive}
\label{p2:homogeneous:primitive}
Fix an integer $r\ge3$, a critical band, and $\varepsilon>0$.
Let $\mathcal P_{r,N}$ be the primitive $r$-element families of original
windows in $I_N$. With the four weights of
Lemma~\ref{p2:homogeneous:envelope},
\begin{equation}
 \sum_{I\in\mathcal P_{r,N}}q_0(I)
       \ll_{r,C,\varepsilon}N^{-r/2+\varepsilon},
 \qquad
 \sum_{I\in\mathcal P_{r,N}}w(I),\quad
 \sum_{I\in\mathcal P_{r,N}}q(I),\quad
 \sum_{I\in\mathcal P_{r,N}}|c(I)|
       \ll_{r,C,\varepsilon}N^{-r/2+\varepsilon}.
 \label{p13:transfer:primitive-bound}
\end{equation}
For each fixed $\eta>0$, the subfamily
\[
 \mathcal P_{r,N}(\eta)
   =\{I\in\mathcal P_{r,N}:\rho_{\mathrm{rel}}(I)\ge\eta L\}
\]
has $O_{r,C,\eta,\varepsilon}(N^\varepsilon)$ members, and each of
its four total weights is $O_{r,C,\eta,\varepsilon}(N^{1-r+\varepsilon})$.
One scale threshold controls these conclusions and every deterministic
subfamily. No regularity, separation, affine compatibility, or positivity
of the original marginals is required. The arity is fixed before that
threshold is chosen.
\end{proposition}

\begin{proof}
The \hyperref[p18:conf:fixed-primitive]{dense-host construction}
in Section~\ref{p14:comp:confinement} proves the required count.
This part of the construction is geometric: it counts centres before
assigning probability weights, and applies to all true windows.
Primitivity makes each relation-space projection injective, whence
\[
 \rho_{\mathrm{rel}}(I)\le\min_{x\in I}d_x\le L.
\]
Lemma~\ref{p2:homogeneous:envelope} therefore gives
\[
 q_0(I)\le(2p)^r2^L\ll_{r,C}N^{1-r}.
\]
Multiplying by the dense host count proves the dense homogeneous bound.
The other three bounds follow pointwise from the same lemma.

For the full homogeneous sum, choose fixed small $\xi,\delta>0$ and
split at $\rho_{\mathrm{rel}}=\xi L$.
Corollary~\ref{p13:half:connected} counts at most
$O_{r,C,\delta}(N^{r/2+\delta})$ primitive hosts.
On the lower part,
$q_0(I)\ll_{r,C}N^{-r+\xi}$, so their contribution is
$O(N^{-r/2+\xi+\delta})$.
On the upper part, the dense bound just proved with parameter $\xi$
gives $O(N^{1-r+\delta})$.
Because $r\ge3$, $1-r<-r/2$.
Choose $\xi,\delta$ within the requested margin and take the maximum
of the finitely many scale thresholds involved. Pointwise domination
again gives the other three sums, and restriction to any deterministic
subfamily preserves every estimate.
\end{proof}

On the regular set, $q_0(I)=p^r2^{\rho_{\mathrm{rel}}(I)}$.
For ordered lists of distinct starts, the same bounds acquire only the
fixed factor $r!$. Estimates at growing cardinalities will instead use
the independent certificate tail below.


The exact centered expansion is
\eqref{p1:alg:centred-identity}. The retained measure $B_2$ uses its
terms of cardinalities zero and two; the stopping identity controls
the remainder through a positive suffix.

\subsection{An independent extension estimate at all cardinalities}

For a fixed number $s$ of source windows and an effective subspace $W$
of their value space, let
\[
 E_W(u)=\sum_{\substack{U\subset\mathcal I_N\setminus\{\text{sources}\}\\
                      W,(S_x)_{x\in U}\text{ not direct}}}
                 u^{|U|}\prod_{x\in U}p_x,
 \qquad u>0.
\]
The weights here are products of marginal probabilities. They are not
conditional probabilities of dependent extensions.

\begin{lemma}[Uniform independent deficit]
\label{p13:transfer:independent-deficit}
For fixed $s,u,C$ and every $\varepsilon>0$,
\begin{equation}
 E_W(u)\ll_{s,u,C,\varepsilon}N^{-1/2+\varepsilon},
 \label{p13:transfer:independent-bound}
\end{equation}
uniformly in the sources and $W$. The sum includes all cardinalities.
\end{lemma}

\begin{proof}
Set $Y=B^2$ and
\[
 a(n)=\prod_{\substack{\ell>Y\\\ell\mid n}}(1+\ell^{-1/2}),\qquad
 A_x=\prod_{\substack{\ell>Y\\\ell\mid n\text{ for some }n\in V_x}}
                                    (1+\ell^{-1/2}).
\]
H\"older applied to the $B$ positions in a window gives
\[
 \sum_{N\leq x<2N}A_x
 \leq\prod_{i=-1}^{L-1}
          \left(\sum_{N\leq x<2N}a(x+i)^B\right)^{1/B}.
\]
Write $a(n)^B=\sum_{d\mid n}h_B(d)$, where $h_B$ is multiplicative,
supported on squarefree integers with primes exceeding $Y$, and
$h_B(\ell)=(1+\ell^{-1/2})^B-1\ll B\ell^{-1/2}$.
The coefficients are nonnegative, so
\[
 \sum_{n\leq3N}a(n)^B
 \leq3N\prod_{Y<\ell\leq3N}\left(1+\frac{h_B(\ell)}\ell\right)
 \leq3N\exp\left(O\left(B\sum_{m>Y}m^{-3/2}\right)\right)
 \ll N.
\]
Thus $\sum_xA_x\ll N$, with a constant independent of $B$.
The small-prime contribution satisfies
\[
 A_{\leq Y}=\prod_{\ell\leq Y}(1+\ell^{-1/2})
       \leq\exp(O(B/\log B))=N^{o(1)}.
\]
Chebyshev's bound and partial summation suffice for this estimate.

Let $P_0$ contain the small primes and the prime factors of the fixed
sources. Their total logarithmic size is $O_s(B\log N)$.
Using the same $Y=B^2$, we have
\[
 \sum_{\substack{\ell\in P_0\\\ell>Y}}\ell^{-1/2}
 \leq\frac{\sum_{\ell\in P_0}\log\ell}{\sqrt Y\log Y}
       =O_s(B/\log B).
\]
The primes up to $Y$ have the same bound by partial summation.
Consequently $A_0=\prod_{\ell\in P_0}(1+\ell^{-1/2})
=\exp(O_s(B/\log B))=N^{o(1)}$.
If $P_k$ also includes the primes of an
arbitrary list $x_1,\ldots,x_k$, then
\[
 \prod_{\ell\in P_k}(1+\ell^{-1/2})
       \leq A_0A_{\leq Y}\prod_{j=1}^k A_{x_j}.
\]
The elementary proof of the exceptional-target estimate
in Proposition~\ref{p1:arith:fixed-branching}, before restricting the
number of sources, gives
\begin{equation}
 D(x_1,\ldots,x_k)
 \ll B(s+k)+B\sqrt N\prod_{\ell\in P_k}(1+\ell^{-1/2}).
 \label{p13:transfer:explicit-exception}
\end{equation}
Contacts account for the first term. Outside them, a selected target
value whose odd prime factors all belong to $P_k$ has the form $a v^2$,
with $a$ squarefree and supported on $P_k$. Summing
$\sqrt{N/a}$ over these $a$ gives the second term, with $B$ possible
placements. This upper bound also includes local defects.

Put $w_x=up_x$, $w_*\ll N^{-1}$,
$M_0=\sum_xw_x=O_{u,C}(1)$, and
$M_1=\sum_xw_xA_x=O_{u,C}(1)$.
Every non-direct extension of size $h$ has an exceptional exposure
in each of its $h!$ orders. Sum over the possible first exceptional
position, allowing repeated centers for an upper bound. The first term
of \eqref{p13:transfer:explicit-exception} contributes at most
\[
 O(Bw_*)\sum_{h\geq1}\frac1{h!}
              \sum_{i=1}^h(s+i-1)M_0^{h-1}=O(B/N).
\]
The second contributes at most
\[
 O(B\sqrt Nw_*A_0A_{\leq Y})
 \sum_{h\geq1}\frac1{h!}\sum_{i=1}^h M_1^{i-1}M_0^{h-i}
       =N^{-1/2+o(1)}.
\]
The double sum is bounded by $e^{\max(1,M_0,M_1)}$.
The varying number of sources has therefore been averaged explicitly;
no constant from a fixed-order theorem has been used at a growing order.
\end{proof}

\begin{lemma}[A fixed number of extensions of a pair]
\label{p13:transfer:fixed-deficit}
For every fixed $k\geq1$,
\begin{equation}
 \sum_{|e|=2}|c(e)|
 \sum_{\substack{U\cap e=\varnothing,\ |U|=k\\
                  S_e,(S_z)_{z\in U}\text{ not direct}}}
                  \prod_{z\in U}p_z
 \ll_{k,C,\varepsilon}N^{-3/2+\varepsilon},
 \qquad S_e=\sum_{i\in e}S_i.
 \label{p13:transfer:fixed-pair-bound}
\end{equation}
\end{lemma}

\begin{proof}
If the new windows are not direct among themselves, choose a primitive
subfamily. A primitive of fixed size $h\geq2$ has independent mass
$O(N^{-h/2+\eta})$ by Theorem~\ref{p13:half:conditional} with no
prescribed sources. The largest such cost is $N^{-1+\eta}$.
After multiplication by $A_2$ and summation of the remaining independent
windows this is $N^{-5/3+O(\eta)}$.

Otherwise choose a minimal nonempty $A\subseteq U$ with
$S_e\cap\sum_{a\in A}S_a\ne0$.
If $|A|=1$, the total cost is the one-target deficit
of Proposition~\ref{p13:activity:deficit}.
If $h=|A|\geq2$, a relation meeting $S_e$ has a nonzero projection in
each of the $h$ target spaces, by minimality. Conditional on $e$,
Theorem~\ref{p13:half:conditional} gives $N^{h/2+\eta}$ possibilities.
The product of the $h$ target marginals costs $O_h(N^{-h})$; summing
the pair weights therefore gives $N^{-2/3-h/2+O(\eta)}$.
The remaining $k-h$ independent weights have bounded sum. There are
only finitely many choices of $h$ and roles because $k$ is fixed.
\end{proof}

\subsection{Uniform certificate tails and the deficit of all pairs}

We record the part of the prime-certificate argument that is uniform in
the number of blocks. A block may be an original window or one of a
fixed finite collection of geometric compressions of width $O(B)$.
Its center ranges over $O(N)$ values and its independent weight is
bounded by a fixed multiple of $p$. A relation of full support selects
a nonzero even vector in every block.

\begin{lemma}[Independent full-support tails]
\label{p13:transfer:independent-tail}
Fix a critical band, fixed bounds on the block widths and centre ranges,
a weight $v>0$, and an exponent $a<2$.
There exist a fixed integer $s\ge3$ and a scale threshold such that,
for every $N$ above that threshold,
\begin{equation}
 \sum_{\substack{|F|\ge s\\F\text{ has a full-support relation}}}
                  \prod_{x\in F}(vp_x)\le N^{-a}.
 \label{p13:transfer:tail-independent}
\end{equation}
The same conclusion holds with one block of any of the fixed compressed
types, weighted by $vp$, and with the remaining factors $vp_x$.
The cutoff and threshold depend only on the fixed parameters, not on
$N,L$ or the chosen blocks. In particular the cutoff is fixed before
any later exponent margin.
\end{lemma}


\begin{proof}
Here are the parameter calculations in the certificate argument of
Lemmas~\ref{p1:core:moderate}--\ref{p1:core:generator}.
Choose $Y=(\log N)^{A_0}$ and a kernel threshold $T=N^b$, with
$A_0>1$ and $0<b<1$. A block containing a value whose odd kernel above
$Y$ is at most $T$ is marked. If
$1/2<\sigma<1$ and $A_0(1-\sigma)<1$, the Rankin count gives
$N^{1-\delta+o(1)}$ marked centers, where
$\delta=(1-b)(1-\sigma)>0$.
Every unmarked block in a full-support relation has two selected
values of kernel exceeding $T$. Their kernels are coprime, since
$Y$ exceeds the fixed block diameter.

For a certificate in which every receiving modulus is at most $O(NB)$,
the CRT and the exponential sum over cardinalities give independent
mass
\[
 N^{-2b+4/A_0+o(1)}.
\]
Here $m=\lceil2\log(O(N))/\log Y\rceil$ and
\[
 \log(B^m(m+1)!)=(4/A_0+o(1))\log N.
\]
If no such certificate exists, each unmarked block has a saturated
neighbor. Its directed degree is $N^{2/A_0+o(1)}$ in both directions,
by Lemma~\ref{p1:core:degrees}.
When $2/A_0+\delta<1$, the component generator of
Lemma~\ref{p1:core:generator}, tilted by $N^{\delta/2}$, stays bounded.
The mass at cardinality at least $s$ is consequently
$N^{-s\delta/2+o(1)}$.
Choose $A_0$ large and then $b$ close enough to one that
$2b-4/A_0>a$; choose an allowed $\sigma$ and then $s$ with
$s\delta/2>a$. All these choices are made before $N$.

The proof uses the actual prime divisors of selected values and a fixed
bound on the block width. Changing one block to a fixed geometric type
changes only constant factors and $B^{O(1)}$ offset choices. Its two
anchors, receiving congruences, and degree bounds obey the same estimates.
Thus the same parameter choices prove the compressed assertion.

The size cutoff can be chosen before any final $\varepsilon$.
Indeed, fix $a<2$, put $a'=(a+2)/2$ and
$\varepsilon'=(2-a)/4$, and apply the established estimate with
these fixed parameters. Since $-a'+\varepsilon'<-a$, it yields the
pure bound $N^{-a}$ above a suitable threshold, for both versions
with the same fixed cutoff. This argument retains $a<2$ and a fixed
bound on the compressed block width.
\end{proof}

For $e=\{x,y\}$ put $E_e(u)=E_{S_e}(u)$ and
$D_2(u)=\sum_{|e|=2}|c(e)|E_e(u)$.

\begin{proposition}[The pair deficit over all cardinalities]
\label{p13:transfer:pair-deficit}
For fixed $u>0$ and every $\varepsilon>0$,
\begin{equation}
 D_2(u)\ll_{u,C,\varepsilon}N^{-3/2+\varepsilon}.
 \label{p13:transfer:pair-deficit-bound}
\end{equation}
In particular it is smaller than the $N^{-4/3}$ scale required below.
\end{proposition}

\begin{proof}
First consider extensions with an internal primitive family of new
windows. At each fixed size their independent mass is $O(N^{-1+\eta})$
by the half-exponent count; at all larger sizes apply
Lemma~\ref{p13:transfer:independent-tail} with any $a>1$.
Summing free complements costs at most $e^{u\mu}$. Their contribution
to $D_2$ is $O(A_2N^{-1+\eta})=O(N^{-5/3+2\eta})$.

For a direct extension, choose its minimal subfamily $A$ crossing $S_e$.
Use Lemma~\ref{p13:activity:pair-partition}. Weak and exceptional pairs
have total activity $O(N^{-1+\eta})$ and $O(N^{-1+\eta})$,
respectively; Lemma~\ref{p13:transfer:independent-deficit} bounds their
complete contribution by $N^{-3/2+O(\eta)}$.
There remain good pairs in a finite set of fixed channels. The cases
$|A|<s$ are covered by Lemma~\ref{p13:transfer:fixed-deficit}.
For $|A|\geq s$, compress the good pair geometrically. Minimality and
directness of $A$ give a full-support relation on the compressed block
and all blocks of $A$. The independent weight assigned to the
compressed center in Lemma~\ref{p13:transfer:independent-tail} is $p$,
whereas the required pair weight is at most $N^{-\kappa_b+o(1)}$,
$\kappa_b=2-1/b\geq5/3$. The ratio is at most
$N^{1-\kappa_b+o(1)}$. Choosing $a=1$ in that lemma bounds this tail
by $N^{-\kappa_b+o(1)}\leq N^{-5/3+o(1)}$.
Each good pair is represented by a fixed-width pattern and one center;
the $B^{O(1)}$ patterns are absorbed in the margins. Free complements
again cost a constant. Choose the finitely many internal margins after
$\varepsilon$ to conclude.
\end{proof}

\subsection{Every fixed connected order reduces to triples}

\begin{proposition}[Reduction of connected activities]
\label{p13:transfer:fixed-activities}
For every fixed $r\geq4$ and every $\varepsilon>0$,
\begin{equation}
 \sum_{\substack{|I|=r\\I\text{ connected}}}|c(I)|
 \ll_{r,C,\varepsilon}N^\varepsilon A_3+N^{-4/3+\varepsilon}.
 \label{p13:transfer:fixed-activity-bound}
\end{equation}
The incompatible part, for every fixed $r\geq4$, is
$O_{r,C,\varepsilon}(N^{-3/2+\varepsilon})$.
\end{proposition}

\begin{proof}
We first bound compatible families by their positive probabilities.
Let $T_3^+$ denote the sum of $q(T)$ over connected compatible triples.
For a primitive triple, $q(T)\leq2c(T)$. For a nonprimitive compatible
triple,
\[
 q(T)=c(T)+\sum_{e\subset T,\ |e|=2}c(e)p_{T\setminus e}
                         +\prod_{i\in T}p_i.
\]
The pair terms sum to the one-target deficit, and the independent term
has mass $O(N^{-3/2+\eta})$ by the half-exponent count. Therefore
\begin{equation}
 T_3^+\leq2A_3+O(N^{-3/2+\eta}).
 \label{p13:transfer:positive-triples}
\end{equation}
Families containing such a triple have mass $N^\eta T_3^+$ after
positive propagation~\eqref{p13:exact:propagation}.
Families containing a primitive of size at least four have mass
$N^{-4/3+\eta}$, by Proposition~\ref{p13:transfer:primitive} and the
same propagation. All seed sizes here are bounded by the fixed $r$.

Consider a remaining compatible connected family. It contains a
dependent pair $e$; otherwise a minimal dependent subfamily would be
one of the seeds just treated. Put $U=S_e$ and choose a nonempty
$A\subset I\setminus e$ minimal with
$U\cap\sum_{a\in A}S_a\ne0$. A single target gives a connected triple,
so $h=|A|\geq2$. If $A$ contains a proper primitive $P$, then
$U\cap\sum_{i\in P}S_i=0$ by minimality. Thus
$q(e\cup P)=q(e)q(P)$. Only $|P|=2$ remains, since larger primitives
were already treated. Compatible dependent pairs satisfy $q(e)\leq2c(e)$,
so the total seed mass is at most $4A_2^2$.
Propagate it to $I$.

Suppose next that $A$ is direct. The nullity $d$ of the compressed
family $U,(S_a)_{a\in A}$ is at most $L$, and
\[
 q(e\cup A)=q(e)\Bigl(\prod_{a\in A}p_a\Bigr)2^d.
\]
For $d<\eta L$, the conditional half-exponent count of the $h$ targets
gives a cost at most $A_2N^{-h/2+O(\eta)}$, hence
$N^{-5/3+O(\eta)}$. For $d\geq\eta L$, fix $h-2$ members of $A$.
The last two targets are each direct from the sum of $U$ and these
fixed targets, and are direct from one another. Otherwise there would
be an internal dependence in $A$ or an earlier crossing of $U$.
Lemma~\ref{p13:transfer:dense-targets} gives conditional weighted cost
$N^{-1+\eta}$. The independent weights of the $h-2$ additional sources
sum to a constant; multiplication by the pair mass gives
$N^{-5/3+O(\eta)}$ again. If $A$ is itself primitive of size at least
three, it was already one of the excluded seeds.

It remains that $A=f=\{y,z\}$ is a dependent pair. By minimality,
$U\cap S_y=U\cap S_z=0$. Put
\[
 d_e=\dim(S_x\cap S_{x'}),\quad d_f=\dim(S_y\cap S_z),\quad
 \delta=\dim(U\cap S_f),\qquad e=\{x,x'\}.
\]
If $\delta<\eta L$, then
$q(e\cup f)=q(e)q(f)2^\delta$, and the total is
$O(N^\eta A_2^2)$.
For $\delta\geq\eta L$, choose $\theta$ from the dense-target lemma.
If $f$ is nonaligned or $d_f\leq\theta L$, that lemma gives
$N^{-5/3+O(\eta)}$ after summing $e$.
If $f$ is a good pair of a fixed channel, fix $e$ and use
Lemma~\ref{p13:transfer:compressed-target}; there are only
$N^\eta$ choices for the compressed center of $f$. Linear algebra gives
\begin{equation}
 d_f+\delta\leq L,
 \qquad q(e\cup f)\leq p_*q(e).
 \label{p13:transfer:two-pair-rank}
\end{equation}
Indeed the projection of $U\cap(S_y+S_z)$ modulo $S_y$ is injective,
and its image has dimension at most $\dim S_z-d_f$.
This case also costs $N^{-5/3+O(\eta)}$.

Every aligned pair with rank greater than $\theta L$ has bounded
channel height, by Lemma~\ref{p1:arith:alignedrank}. Outside the good
pairs at those finitely many heights, the normalization estimates in
Theorem~\ref{p1:arith:pairs} and
 Lemma~\ref{p1:arith:fixed-channel-mass} give at most
$N^{1/2+\eta}$ exceptional aligned pairs. The endpoint pairs of height
two and the pairs meeting an irregular window satisfy the same bound;
height one has rank $o(L)$ when compatible.
Call this exceptional collection $\mathcal B$.
If only $f$ belongs to $\mathcal B$, exchange $e$ and $f$ and use one
of the preceding cases. This exchange is valid because a target in
$e$ intersecting $S_f$ would form a connected triple, already removed.
If both belong to $\mathcal B$, there are at most $N^{1+2\eta}$ hosts.
The symmetric forms of \eqref{p13:transfer:two-pair-rank} give
\[
 d_e+d_f+\delta\leq L+\min(d_e,d_f)\leq\tfrac32L+o(L),
\]
where the final inequality uses
Theorems~\ref{p13:rig:nonaligned} and~\ref{p13:rig:aligned}.
The joint probability is at most $N^{-5/2+o(1)}$, so this contribution
is $N^{-3/2+O(\eta)}$. These cases exhaust the possible minimal sets
$A$. Positive propagation adds the remaining windows, proving the
compatible part of \eqref{p13:transfer:fixed-activity-bound}.

For incompatible families we use induction on $r$, starting with the
triple estimate. The exact identity is
\begin{equation}
 c(I)=-\prod_{i\in I}p_i-
          \sum_{\substack{J\subsetneq I\\|J|\geq2}}
                     c(J)\prod_{i\in I\setminus J}p_i.
 \label{p13:transfer:incompatible-expansion}
\end{equation}
The independent term over connected $r$-tuples is
$O(N^{-r/2+\eta})$. The pair terms cost $O(N^{-3/2+\eta})$ by
Lemma~\ref{p13:transfer:fixed-deficit}.
For a fixed proper $J$ of size at least three, connectedness of $I$
requires a non-direct extension of its total space; its independent
cost is $N^{-1/2+\eta}$ by
Lemma~\ref{p13:transfer:independent-deficit}.
Canonical factorization of $c(J)$ has no singleton factors.
By induction a factor of size at least three costs
$N^\eta A_3+N^{-4/3+\eta}$; a decomposition consisting only of pairs
has at least two of them and costs $O(A_2^2)$.
There are finitely many partitions at each fixed size. Multiplying by
the independent extension cost proves
\eqref{p13:transfer:fixed-activity-bound} at order $r$.
Inserting the arithmetic bound $A_3\ll N^{-41/30+\eta}$ from
Corollary~\ref{p13:third:activity} into the same induction gives the
stated $N^{-3/2+\varepsilon}$ incompatible activity. No law
approximation or lower bound has been used in this induction.
\end{proof}

\subsection{A uniform tail after 256 blocks}

The preceding proposition has constants depending on its fixed order.
We now use the certificate generator to handle all larger orders in
the stopped expansion. The number $256$ is a sufficient fixed cutoff,
not an optimal threshold.

\begin{lemma}[Certificate budgets with several anchors]
\label{p13:transfer:anchor-budgets}
Choose an even integer $t\geq2$, $A_0>1$, $0<b<1$, and
$1/2<\sigma<1$ with $A_0(1-\sigma)<1$. Put
\[
 \delta=(1-b)(1-\sigma),\qquad \theta=t/A_0.
\]
Suppose each unmarked block in a family has a relation selecting at
least $t$ values in that block; the relation may depend on the block.
Mark a block when some value has odd kernel above
$Y=(\log N)^{A_0}$ at most $N^b$.
Under any fixed independent weight per block, the moderate certificates
have total mass
\begin{equation}
 N^{-tb+2t/A_0+o(1)}.
 \label{p13:transfer:moderate-budget}
\end{equation}
If $\theta+\delta<1$, the saturated families of size at least $s$
have total independent mass
\begin{equation}
 N^{-s\delta/2+o(1)}.
 \label{p13:transfer:saturated-budget}
\end{equation}
\end{lemma}

\begin{proof}
This is the calculation in Lemma~\ref{p13:transfer:independent-tail}
with $t$ anchors. Their pairwise coprime kernels have product exceeding
$N^{tb}$ and at most
$m_t=\lceil t\log(O(N))/\log Y\rceil$ prime factors.
The CRT bound and exponential cardinality sum have combinatorial cost
$B^{m_t}(m_t+1)!=N^{2t/A_0+o(1)}$, proving
\eqref{p13:transfer:moderate-budget}.
The saturated graph has degree $N^{\theta+o(1)}$ in both directions;
the number of marked vertices is $N^{1-\delta+o(1)}$.
After tilting by $N^{\delta/2}$, the three terms of the component
generator of Lemma~\ref{p1:core:generator} are
\[
 N^{-\delta/2+o(1)},\quad
 N^{-1+\theta+\delta/2+o(1)},\quad
 N^{-1+\theta+\delta+o(1)}.
\]
They remain bounded, and in fact tend to zero, under the stated
condition. Removing the tilt proves
\eqref{p13:transfer:saturated-budget}. The inequality $p_x\leq2p$
is absorbed by doubling the fixed weight. The arguments concern true
prime factors and permit contacts and irregular source windows.
\end{proof}

For a family $D$, let $\mathcal K(D)$ be the kernel of the sum map
$\bigoplus_{x\in D}S_x\to\sum_{x\in D}S_x$, and set
\[
 d_x(D)=\dim\operatorname{proj}_{S_x}\mathcal K(D)
       =\rho(D)-\rho(D\setminus\{x\}).
\]
Call $D$ a $10$-core when each $d_x(D)\geq9$.
The effective projection admits a linear section in the formal even
coordinates. Setting nine pivot coordinates to one gives a word of
weight at least nine; even parity makes its weight at least ten.
Thus the preceding budgets apply with $t=10$.
For a compatible family Fourier expansion gives exactly
\begin{equation}
 \frac{c(D)}{q(D)}
 =\P_{v\ \mathrm{uniform\ in}\ \mathcal K(D)}
       (\operatorname{proj}_{S_x}v\ne0\text{ for every }x\in D)
 \geq1-|D|2^{-9}.
 \label{p13:transfer:core-ratio}
\end{equation}
In particular $q(D)\leq2c(D)$ for $|D|<256$.

\begin{proposition}[Uniform connected and primitive tails]
\label{p13:transfer:uniform-tail}
For fixed $w>0$, every $\varepsilon>0$, and every $K_N=o(L)$,
\begin{align}
 \sum_{\substack{I\text{ connected},\ 256<|I|\leq K_N\\
                         \rho_*(I)\leq2L}}
              w^{|I|}|c(I)|&\ll N^{-4/3+\varepsilon},
                    \label{p13:transfer:connected-tail}\\
 \sum_{\substack{P\text{ primitive}\\|P|>256}}
              w^{|P|}|c(P)|&\ll N^{-4/3+\varepsilon}.
                    \label{p13:transfer:primitive-tail}
\end{align}
The second sum includes all cardinalities. Here
$\rho_*(I)=\max\{\rho(A):A\subseteq I,\ q(A)>0\}$.
\end{proposition}

\begin{proof}
For compatible $10$-cores of nullity at most $2L$, choose
\[
 t=10,\quad A_0=12,\quad b=31/60,\quad
 \sigma=23/25,\quad\delta=29/750,\quad\theta=5/6.
\]
The condition $\theta+\delta<1$ is strict. Multiplying the independent
budgets by $2^{2L}=O_C(N^2)$ gives, at sizes at least $256$,
\begin{equation}
 \sum_D w^{|D|}q(D)
 \ll N^{-3/2+o(1)}+N^{-1106/375+o(1)}.
 \label{p13:transfer:large-core}
\end{equation}
Indeed $2-10(31/60)+20/12=-3/2$ and
$2-256(29/1500)=-1106/375$.

If $\rho_*(I)\leq L/60$, the centered expansion gives
\[
 |c(I)|\leq2^{|I|}\Bigl(\prod_{x\in I}p_x\Bigr)2^{L/60}.
\]
A nonzero centered coefficient has a full-support Fourier relation,
so two anchors are available in every block. Use
\[
 t=2,\quad A_0=12,\quad b=17/20,\quad
 \sigma=23/25,\quad\delta=3/250,\quad\theta=1/6.
\]
Starting at cardinality $256$, the two costs after the rank factor are
\begin{equation}
 N^{-27/20+o(1)}+N^{-2279/1500+o(1)}.
 \label{p13:transfer:low-rank-tail}
\end{equation}
The four costs in \eqref{p13:transfer:large-core} and
\eqref{p13:transfer:low-rank-tail} leave the exact margins
\[
 \begin{aligned}
 \frac32-\frac43&=\frac16,&
 \frac{1106}{375}-\frac43&=\frac{202}{125},\\
 \frac{27}{20}-\frac43&=\frac1{60},&
 \frac{2279}{1500}-\frac43&=\frac{93}{500}.
 \end{aligned}
\]
Thus the smallest margin is $1/60$. For the fixed certificate parameters
and fixed weight $w$, every displayed $o(1)$ loss is less than $1/120$
for sufficiently large $N$; no certificate parameter is allowed to vary
with $N$.

For the remaining connected $I$, choose deterministically a compatible
$A\subseteq I$ attaining $\rho_*(I)>L/60$. Delete blocks of projection
dimension at most eight until a $10$-core $D$ remains. Since
$8K_N=o(L)$, its nullity exceeds $L/120$ for large $N$, so $|D|\geq2$.
With $r=|I|$ and $h=|D|$, the centered expansion and maximality of
$\rho_*(I)$ imply
\begin{equation}
 |c(I)|\leq2^r2^{8(r-h)}
               \Bigl(\prod_{x\in I\setminus D}p_x\Bigr)q(D).
 \label{p13:transfer:peeling}
\end{equation}
The elements outside $A$ are already in this product; there is no
enumeration of deletion orders. All fixed powers of two can be absorbed
in the independent weight of the complements.

If $h\geq256$, sum the complements freely and apply
\eqref{p13:transfer:large-core}. If $3\leq h<256$, then
$q(D)\leq2c(D)$ by \eqref{p13:transfer:core-ratio}.
Proposition~\ref{p13:transfer:fixed-activities}, the arithmetic triple
bound, and canonical factorization show that the sum of $|c(D)|$ at
each such fixed size is $O(N^{-1+\eta})$. Because $D$ is proper in
the connected $I$, its complement is a non-direct independent extension.
Lemma~\ref{p13:transfer:independent-deficit} adds $N^{-1/2+\eta}$.
If $h=2$, then $D=e$ is a compatible dependent pair and $q(e)\leq2c(e)$;
the complement is exactly the pair deficit of
Proposition~\ref{p13:transfer:pair-deficit}. These bounds prove
\eqref{p13:transfer:connected-tail}.

A primitive family has nullity at most $L$. If it is incompatible,
$|c(P)|=\prod_Pp_x$ and the independent two-anchor budgets suffice.
If it is compatible with nullity at most $L/60$, use
\eqref{p13:transfer:low-rank-tail}. Otherwise each projection has
dimension $\rho(P)>L/60$, so it is itself a $10$-core for large $N$.
Apply \eqref{p13:transfer:large-core} directly, without an upper
cardinality cutoff. This proves \eqref{p13:transfer:primitive-tail}.
\end{proof}


\subsection{Conditional dense targets and small-kernel counts}
\label{p14:comp:transfer}

We prove the conditional dense-pair estimate of
Lemma~\ref{p13:transfer:dense-targets} and then the separate auxiliary
small-kernel estimate of Lemma~\ref{p13:transfer:small-kernels}.
The first supplies a fixed-source input to the reconstruction;
the second completes the auxiliary statement recorded above.
All source counts, widths, and rank thresholds are fixed before
$N$ tends to infinity. Neither proof uses the finite-transfer theorem
or the law approximation.

\subsubsection{Proof of Lemma~\ref{p13:transfer:dense-targets}}
\label{p14:comp:transfer-dense-targets}

\begin{proof}
We first count target pairs, keeping the number of binary components in
\emph{each} target explicit. All component counts below include isolated
active coordinates.

\paragraph{The two injective projections.}
Put
\[
 R=\ker\bigl(W\oplus S_y\oplus S_z\longrightarrow W+S_y+S_z\bigr),
 \qquad \dim R=d.
\]
If an element of $R$ has zero $y$-projection, its other two entries
belong to $W\cap S_z=0$. Hence the projection $R\to S_y$ is
injective; the other projection is injective by $W\cap S_y=0$.
In particular $d\leq\min(\dim S_y,\dim S_z)\leq L$.
Represent $W$ by a linear section in the physical source coordinates,
and each target space by a section in its formal even-selection space.
These sections preserve both injective target projections and all actual
prime-valuation equations, including on an irregular window.

Fix $0<\delta<1/2$, to be chosen inside the final exponent margin,
and choose fixed positive $\xi,\theta$ with
\[
 \xi\leq\tfrac14\min(\eta,\delta),\qquad
 \theta\leq\tfrac18\min(\eta,\delta).
\]
Apply Lemma~\ref{p1:fixed:recovery} to obtain $R'\subseteq R$ with
$d':=\dim R'\geq d-\xi L$. Choose its recovery parameter large
enough to give this loss and to control circuits in the $s+2$ physical
blocks. Every nonzero coordinate label on $R'$ then has a fixed
bounded multiplicity $K$.

For each prime $\ell>H\asymp B$, pair equal active labels in its
row. All edges of the resulting binary graph join equal nonzero
labels, so every component lies within one label class and has at
most $K$ vertices. The coordinate labels in either target span
$(R')^*$, since its projection is injective. A component contributes
only one label to this span, even if it meets that target several times.
Consequently each target meets at least $d'$ binary components.
This is a separate lower bound for each target, not just for their union.

\paragraph{Enumerating coefficients from the fixed sources.}
After the equal labels have been paired, the remaining distinct labels
in a prime row split into minimal circuits. Each such circuit uses at
least three physical blocks: $\ell>H$ has at most one occurrence in
any block. Since only two blocks are free, every circuit contains a
source occurrence. The product of primes assigned to a circuit support
therefore divides a fixed source value.
The recovery bound gives a fixed bound on the number of circuit supports
incident to any selected coordinate. For a fixed number of selected
binary components, enumerate their positions, these circuit supports,
one source occurrence on each support, and divisors of the corresponding
source values. The cost is
$B^{O(1)}N^{o(1)}=N^{o(1)}$. Small-prime parities cost
$2^{O(\pi(H))}=N^{o(1)}$ as well. Thus the squarefree coefficients
of their component equations have $N^{o(1)}$ choices and height
$N^{A_0}$ for a fixed $A_0$. The height bound per component is fixed
once $s$, the recovery parameters and the width multiplier are fixed;
it does not increase when more disjoint component equations are selected.
The number of equations required by Lemma~\ref{p1:arith:geometric}
is chosen afterwards. No subspace or full deletion pattern is enumerated.

\paragraph{Removing the small aligned curves.}
Two exact units between the targets determine a reduced alignment of
height $O(B)$. Two distinct alignments of this height would force
$y,z=O(B^3)$, impossible in the critical range. If there are fewer than
two exact units, no removal is needed. Otherwise let $u\geq2$ be the
number of units on the unique alignment. Their even selections have dimension
$u-1$ before passage to square classes, and a target has at most one
local rank defect. Their common image therefore gives
\[
 u-2\leq\dim(S_y\cap S_z)\leq\theta L
\]
in the weak aligned case. Keep at most one of these units and exclude
from selection every binary component meeting an endpoint of another
unit. At most $2(u-1)\leq2\theta L+2$ components are excluded.
This deletes components from the selectable collection; it does not
restrict $R'$ again. The surviving positions contain no two exact units
on an alignment, which excludes the small-curve alternatives in the
proof of Lemma~\ref{p1:arith:geometric}. No removal is needed for a
nonaligned pair.

For $t=y,z$, write $m_t$ for the number of selectable components
meeting target $t$. The separate projection bounds now give
\begin{equation*}
 m_t\geq d'-2\theta L-2
      \geq d-(\xi+2\theta)L-2.
\end{equation*}
Thus $m_t\geq\eta L/3$ for large $N$ whenever $d\geq\eta L$;
if $d>(1/2+\delta)L$, then
$m_t>(1/2+\delta/3)L$ for both targets.

\paragraph{The first stratum: a positive rank proportion.}
A component type records its number of coordinates in each physical
block. There are finitely many types, because the component size is at
most $K$. If a type meeting both targets occurs sufficiently often,
select the fixed number of its components required by the two-variable
geometric count. With the small curves excluded above and coefficients
of height $N^{A_0}$, the proof of Lemma~\ref{p1:arith:geometric}
gives $N^{1/2+\varepsilon'}$ target pairs, including the
$N^{o(1)}$ pattern and coefficient choices.

Otherwise only a bounded number of components meet both targets.
Each target still has a positive proportion of $L$ components avoiding
the other target. In each target choose a univariate type occurring with
positive frequency, and then the fixed number of its components needed
by the one-variable part of Lemma~\ref{p1:arith:geometric}, taking
at least two. Their component equations form a one-variable system
with fixed coefficients and at least two distinct selected offsets;
different components have disjoint vertices. In the degree-one case,
these are precisely the two linear factors for
Lemma~\ref{p1:arith:pell}. The one-variable count gives
$N^{\varepsilon'}$ target pairs after allocating the margin between
the two centers and the coefficient choices. Every selected set has fixed size, so its even-selection
square-class map is injective for large $N$ by
Lemma~\ref{p1:arith:runge}; irregularity of the entire window is
irrelevant to this local use. We have proved
\[
 \#\{(y,z):d\geq\eta L\}\ll N^{1/2+\varepsilon'}.
\]
Both counts here and below impose the pair restriction in the lemma.

\paragraph{The second stratum: rank greater than half a window.}
Suppose $d>(1/2+\delta)L$. Let $n_{1,t}$ be the number of selectable
components using exactly one coordinate in target $t$. Every other
component meeting that target uses at least two of its $B=L+1$
coordinates. Hence, separately for $t=y,z$,
\[
 B\geq n_{1,t}+2(m_t-n_{1,t}),\qquad
 n_{1,t}\geq2m_t-B>\frac{2\delta}{3}L-1.
\]
If a mixed type with one coordinate in either target occurs sufficiently
often, the degree-one case of the geometric count gives
$N^{\varepsilon'}$ pairs directly. Otherwise only boundedly many of
the $n_{1,t}$ components are mixed. The displayed inequality then
provides two components for each target, each using one coordinate of
that target and no coordinate of the other. The resulting two distinct
linear factors give the uniform two-value Pell count for each center.
Thus
\[
 \#\{(y,z):d>(1/2+\delta)L\}\ll N^{\varepsilon'}.
\]
This explains the threshold $L/2$: above it, the coordinate budget
forces singleton occurrences in each target.

Finally $p_y,p_z\ll N^{-1}$ and $2^L\asymp_C N$.
The stratum $\eta L\leq d\leq(1/2+\delta)L$ has weighted cost
$N^{1/2+\varepsilon'}N^{-2+1/2+\delta}$, and the stratum above it
has cost $N^{\varepsilon'}N^{-2+1}$. Choose $\delta>0$ and then
$\varepsilon'>0$ so that $\delta+\varepsilon'<\varepsilon$;
all preceding choices are fixed before $N$. This proves
\eqref{p13:transfer:dense-target-bound}.
In a primitive family, each target is direct from the prescribed
$r-2$ sources separately, as well as from the other target; the full
nullity is at most $L$. Multiplication by the $r-2$ source marginals
and summation over the prescribed source collection gives
\eqref{p13:transfer:restricted-dense}, after dividing the final exponent
margin between the source count and the conditional estimate.
\end{proof}

\subsubsection{Proof of Lemma~\ref{p13:transfer:small-kernels}}
\label{p14:comp:transfer-small-kernels}

\begin{proof}
For the first assertion fix the $m$ offsets, at a polylogarithmic cost,
and write $d_i=n_i/K_Y(n_i)=s_i a_i^2$. Since
$d_i\gg N^{1-b}$ and $\gcd(d_i,d_j)\ll B$,
\[
 \operatorname{lcm}(d_1,\ldots,d_m)
 \geq\frac{\prod_i d_i}{\prod_{i<j}\gcd(d_i,d_j)}
 \geq N^{m(1-b)-o(1)}>N.
\]
The $a_i$ are pairwise coprime. If $\prod_i a_i^2>N$, multiply the
square factors in a fixed order until their product first exceeds $N$.
The resulting square $D$ lies between $N$ and $N^{1+u}$, so it has
at most $N^{(1+u)/2}$ possible roots. Assign its complete prime powers
to the $m$ offsets; the number of assignments is
$m^{\omega(D)}=N^{o(1)}$. The CRT leaves at most one center modulo
$D>N$ in the interval.

If $\prod_i a_i^2\leq N$, start with this product and add the small
prime factors from the least common multiple until the first crossing
of $N$. Now $N<D\leq NY$ and $D=sa^2$ with $s$ $Y$-smooth.
For fixed $\sigma>\max(1/2,1-1/A)$, the number of such modules is at most
\[
 (NY)^\sigma\sum_{a\geq1}a^{-2\sigma}
       \prod_{\ell\leq Y}(1-\ell^{-\sigma})^{-1}
       =N^{\sigma+o(1)}.
\]
Choose $\sigma$ just above $1-1/A$. Assigning the complete prime powers
again costs $N^{o(1)}$ and the CRT gives at most one center. This proves
\eqref{p13:transfer:small-kernels-bound}.

For the second assertion, use $H=C_DB$ as in
Lemma~\ref{p13:half:small-kernels}.
That lemma gives the result for $q>0$.
For $q=0$, use $\mathcal X_0\subseteq\mathcal X_{\varepsilon/2}$
and apply the lemma with error margin $\varepsilon/2$.
\end{proof}
